\section{Results and discussion}
\label{sec:results}

All autoencoder numbers in this section are validation results.  The test
split has not been read.

\subsection{PCA: the linear compression baseline}

PCA finds orthogonal linear patterns that account for the largest variation in
the training surfaces.  Each surface is flattened from
$31\times101$ to 3,131 values, centered using the training mean, and projected
onto the leading patterns.  The first two PCA components retain about 96.23\%
of validation variance.  On the same validation split, ten components recover
about 99.94\% of the variance.  Separately, the rank-25 PCA fit accounts for
99.9986\% of training variance.

Figure~\ref{fig:pca-variance} shows the individual and cumulative training
variance, while Figure~\ref{fig:pca-rank} shows validation reconstruction
error as the rank increases.  Most of the gain appears in the first few
components.

\reportfigure{pca_explained_variance.pdf}{Individual and cumulative variance
explained by PCA fitted to the training surfaces.  Table~\ref{tab:preliminary-results}
instead reports reconstruction on the validation set.}{fig:pca-variance}

\reportfigure{pca_reconstruction_error_by_rank.pdf}{Validation reconstruction
error as the number of PCA components increases.  Rank 2 is the matched
baseline for the autoencoder; rank 10 is shown as a higher-dimensional
reference.}{fig:pca-rank}

\subsection{Autoencoder result with a two-number latent representation}

The autoencoder measurements currently available come from an earlier
25-epoch trial, not from the longer training configuration described in
Section~\ref{sec:autoencoder}.  Its best checkpoint, at epoch 20, recovers
98.09\% of the validation variance with a two-number latent space.  The
validation RMSE is 0.03795 in $\log_{10}\mathcal A_\Theta$, the mean relative
error in physical amplitude is 6.48\%, and the 95th percentile of per-surface
relative RMSE is 13.95\%.  Only the scalar summaries from that trial are
available in this report; its checkpoint and training history are not present
in the current repository.  The comparison is therefore preliminary.

\subsection{A fair comparison at equal latent dimension}

Table~\ref{tab:preliminary-results} compares methods on the same validation
indices and the same log-amplitude target.  The main comparison is autoencoder
dimension 2 versus PCA rank 2.  PCA rank 10 is included only as a
higher-dimensional accuracy reference.

\begin{table}[t]
\centering
\caption{Preliminary validation comparison on the stored 10,000-surface
validation
split.  The autoencoder uses a two-number bottleneck.  PCA rank 2 is the
matched-dimension baseline.  PCA rank 10 is a higher-dimensional reference,
not a competing two-number method.}
\label{tab:preliminary-results}
\begin{tabular}{@{}lrrrr@{}}
\toprule
Method & Latent size & Variance recovered & Log RMSE & Mean relative error \\
\midrule
PCA & 2 & 0.9623 & 0.0533 & 0.0866 \\
Conv1D autoencoder & 2 & 0.9809 & 0.0380 & 0.0648 \\
PCA & 10 & 0.9994 & 0.0068 & 0.0104 \\
\bottomrule
\end{tabular}
\end{table}

In this trial, the autoencoder has higher recovered variance and lower
reconstruction errors than PCA at the same latent dimension.  This suggests
that two linear components do not capture the model family as efficiently as
the nonlinear mapping.  It does not show that two numbers are sufficient:
the 99.9\% variance target was not reached, and the 95th-percentile result
shows that some surfaces are substantially harder to reconstruct.

\subsection{Latent-space interpretation}

We have not yet mapped the two latent coordinates against the 13 input
parameters, so neither coordinate is given a physical interpretation.  Such an
analysis should examine both parameter correlations and reconstructed surfaces
along paths through the latent space.

\subsection{Limitations}

These results have four main limitations:

\begin{itemize}
  \item The analytic response is a deliberately flexible, phenomenological
        model family.  It has not yet been calibrated to a particular survey or
        shown to describe all physical intrinsic-alignment effects.
  \item The autoencoder numbers above come from validation data.  The untouched
        test split will be evaluated only after the architecture and latent
        dimension have been chosen.
  \item The autoencoder result comes from one short run whose checkpoint and
        training history are unavailable; it has not been repeated across
        several random seeds.
  \item Reconstruction accuracy varies across surfaces, so the mean metrics do
        not establish adequate precision at every scale, redshift, or
        parameter-space location.
\end{itemize}

\subsection{Next steps}

\begin{enumerate}
  \item compare latent dimensions 2, 4, 6, 8,
        and 10, using the same data and metrics;
  \item repeat the most promising model with several random seeds;
  \item map reconstruction error over $(k,z)$ and over the 13-parameter space;
  \item test whether a physically motivated weighted loss improves important
        regions without worsening the rest of the surface; and
  \item after selecting the method, evaluate the untouched test set.
\end{enumerate}
